% Annale FICTIVE (aucun contenu réel) : examen britannique à essais, numéro de candidat, notes en marge.
\documentclass[12pt,a4paper]{article}
\usepackage[T1]{fontenc}
\usepackage[utf8]{inputenc}
\usepackage{mathptmx}
\usepackage[margin=2.5cm]{geometry}
\usepackage{fancyhdr}
\pagestyle{fancy}\fancyhf{}
\renewcommand{\headrulewidth}{0pt}
\fancyfoot[C]{\thepage}
\fancyhead[R]{\small IMAG2001}
\newcommand{\qq}[2]{\par\bigskip\noindent\textbf{#1.}\quad #2\hfill\textbf{[25 marks]}\par}
\begin{document}
\thispagestyle{empty}
\begin{center}
{\large UNIVERSITY OF NOWHERE}\\[10pt]
{\Large\textbf{IMAG2001 Imaginary Literatures}}\\[6pt]
Summer Examination\\[4pt]
Time allowed: TWO hours
\end{center}
\vspace{10pt}
\noindent Candidate number: \makebox[0.35\textwidth]{\hrulefill}\par\vspace{6pt}
\noindent Desk number: \makebox[0.2\textwidth]{\hrulefill}\par\vspace{12pt}
\noindent\textbf{INSTRUCTIONS TO CANDIDATES}
\begin{enumerate}\setlength{\itemsep}{2pt}
\item Answer TWO questions, one from each section.
\item Write your answers in the answer booklet provided.
\item Do not turn over this page until instructed to do so.
\end{enumerate}
\clearpage
\noindent\textbf{SECTION A}
\qq{1}{``An invented novel is always a map of a place that never existed.'' Discuss with reference to at least two fictional works.}
\qq{2}{Assess the claim that imaginary poetry is defined by its silences.}
\par\bigskip\noindent\textbf{SECTION B}
\qq{3}{To what extent can a made-up chronicle be trusted as a historical source?}
\qq{4}{Compare the narrators of two invented epics. How do they shape the reader's sympathy?}
\end{document}
